% =====================================================================
% MEJORADO (27 oct 2026): integración #5-#7, métricas de código, verificación y suites de pruebas.
\subsection{Fase de estabilización}
\label{sec:des-estabilizacion}

En esta fase se concentró el trabajo en integrar los módulos, uniformar la
arquitectura y verificar la calidad del código. Las adiciones funcionales
corresponden al componente de servidor (funciones \archivo{buscarAlimentos} y
\archivo{guardarRutina}), que se incorporó para proteger la clave de la API de
nutrición y para validar en el servidor la creación de rutinas, y, el 3 de
noviembre de 2026, al reconocimiento de máquinas en la nube
(\archivo{reconocerMaquina}) y al módulo nativo de ARCore
(sección~\ref{sec:des-produccion}).

\subsubsection{Integración de los módulos}

Los módulos desarrollados en las ramas \archivo{feature/auth},
\archivo{feature/machines} y \archivo{feature/ejercicios} se integraron en
\archivo{develop} mediante las solicitudes de integración \#1, \#3 y \#4. Los
módulos de rutinas, escáner, nutrición y progreso se desarrollaron en
\archivo{feature/routines}, integrada mediante la solicitud \#5 (18 de octubre
de 2026). La solicitud \#6 (19 de octubre) incorporó el cambio de nombre de los
archivos a nombres en inglés (\archivo{refactor/english-file-names}) y la
solicitud \#7 (27 de octubre) integró la rama \archivo{feature/modelAI}, con los
modelos de TensorFlow Lite, el componente de servidor y las pruebas de
seguridad. La navegación principal quedó unificada en la barra
de pestañas de \archivo{(app)} (Inicio, Máquinas, Progreso y Nutrición), con
las rutinas y el perfil accesibles desde el inicio, y la protección de rutas
centralizada en el diseño raíz.

\subsubsection{Refactorización por capas}

Durante las primeras iteraciones, solo el módulo de nutrición cumplía
completamente la separación en capas; las pantallas de autenticación, de
máquinas y de rutinas llamaban a Firebase directamente. En la estabilización
se extendió el patrón a todo el sistema, de modo que cada módulo cuenta con su
controlador, sus entidades, su esquema de validación y su repositorio
(Tabla~\ref{tab:modulos-capas}).

\begin{table}[H]
\centering
\caption{Componentes de cada módulo por capa}
\label{tab:modulos-capas}
\small
\begin{tabularx}{\textwidth}{|L{2.4cm}|L{4.5cm}|Y|L{2.7cm}|}
\hline
\rowcolor{encabezado} \textbf{Módulo} & \textbf{Controlador} & \textbf{Modelo (entidad /
repositorio)} & \textbf{Vista} \\
\hline
Autenticación & \archivo{AuthController} & \archivo{Usuario} /
\archivo{userRepository} & \archivo{(auth)}, \archivo{(onboarding)},
\archivo{profile} \\ \hline
Máquinas & \archivo{MachinesController} & \archivo{Maquina} /
\archivo{machineRepository} & \archivo{machines/*} \\ \hline
Ejercicios & \archivo{ExercisesController}, \archivo{WgerExercisesController} &
\archivo{Ejercicio}, \archivo{EjercicioWger} / repositorios de ejercicios &
Modales de selección \\ \hline
Rutinas & \archivo{RoutinesController} & \archivo{Rutina} /
\archivo{routineRepository} & \archivo{routines/*} \\ \hline
Nutrición & \archivo{NutritionController} & \archivo{RegistroNutricional},
\archivo{AlimentoUSDA} / \archivo{nutritionRepository},
\archivo{foodRepository} & \archivo{nutrition} \\ \hline
Progreso & \archivo{ProgressController} & \archivo{ProgresoFisico} /
\archivo{progressRepository} & \archivo{progress/*}, inicio \\ \hline
Entrenamientos & \archivo{WorkoutsController} & \archivo{Entrenamiento}
/ \archivo{workoutRepository} & \archivo{progress/*} \\
\hline
\end{tabularx}
\end{table}

Otras mejoras aplicadas en la refactorización fueron:

\begin{itemize}
    \item \textbf{Tiempo real en lugar de recarga manual.} Las pantallas que
    consultaban los datos una sola vez y ofrecían una función de recarga
    pasaron a suscribirse con \archivo{onSnapshot} a través de los métodos
    \archivo{observar} de cada controlador, con la cancelación de la
    suscripción al desmontar la pantalla.
    \item \textbf{Manejo de errores de las suscripciones.} Se añadió un
    manejador por defecto que, ante un error de lectura, registra un aviso y
    entrega un resultado vacío para que la interfaz no quede esperando
    indefinidamente. El error de permisos que se produce al cerrar sesión con
    suscripciones activas se reconoce como esperado y se ignora.
    \item \textbf{Limpieza de campos indefinidos.} Firestore rechaza los
    campos con valor \archivo{undefined}; en lugar de exigir que cada
    repositorio lo tenga en cuenta, la limpieza se centralizó en
    \archivo{firestoreService}, el único archivo que usa el SDK.
    \item \textbf{Validación en el controlador.} Todas las operaciones de
    escritura validan sus datos con Zod antes de llegar al repositorio, y las
    reglas de seguridad repiten las validaciones críticas en el servidor.
    \item \textbf{Reglas alineadas con el modelo real.} Las reglas de
    Firestore se revisaron campo por campo contra las entidades y los
    repositorios, de modo que cada escritura que realiza la aplicación
    satisface su regla y ninguna otra es aceptada.
    \item \textbf{Convenciones de código.} Importaciones con el alias
    \archivo{@/} en lugar de rutas relativas largas, formato uniforme con
    Prettier y reglas de análisis estático de React (incluidas las de
    \emph{hooks}) sin desactivar.
\end{itemize}

\subsubsection{Métricas del código}

La Tabla~\ref{tab:metricas} resume el tamaño del código fuente por capa al
cierre de la estabilización.

\begin{table}[H]
\centering
\caption{Tamaño del código fuente por capa}
\label{tab:metricas}
\begin{tabular}{|L{3.0cm}|L{5.4cm}|C{1.8cm}|C{1.8cm}|}
\hline
\rowcolor{encabezado} \textbf{Capa} & \textbf{Carpetas} & \textbf{Archivos} & \textbf{Líneas} \\
\hline
Vista & \archivo{app/}, \archivo{components/} & 87 & 11\,410 \\ \hline
Controlador & \archivo{controllers/}, \archivo{utils/}, \archivo{store/},
\archivo{hooks/} & 41 & 5\,454 \\ \hline
Modelo & \archivo{models/}, \archivo{schemas/} & 32 & 1\,852 \\ \hline
Servicios & \archivo{services/}, \archivo{config/} & 20 & 1\,423 \\ \hline
Otros & \archivo{constants/}, \archivo{types/} & 6 & 93 \\
\hline
\textbf{Total} & & \textbf{186} & \textbf{20\,232} \\
\hline
\end{tabular}
\fuente{Conteo de archivos \archivo{.ts} y \archivo{.tsx} de la carpeta
\archivo{src/}, sin las pruebas automáticas, incluidos comentarios y líneas en
blanco, al 3 de noviembre de 2026. No incluye el código de las funciones de
servidor (\archivo{functions/}) ni el del módulo nativo de ARCore
(\archivo{modules/coach-ar/}, 606 líneas de Kotlin en cuatro archivos). El
crecimiento respecto del 27 de octubre se debe sobre todo a los componentes
compartidos de interfaz, a los servicios y esquemas del reconocimiento en la
nube y a la interfaz tipada de la vista de realidad aumentada.}
\end{table}

\subsubsection{Verificación estática}

Como criterio de cierre de la fase se estableció una secuencia de
verificaciones automáticas que el código debe superar
(Tabla~\ref{tab:verificacion-estatica}).

\begin{table}[H]
\centering
\caption{Verificaciones estáticas del código}
\label{tab:verificacion-estatica}
\begin{tabularx}{\textwidth}{|L{4.0cm}|Y|L{3.4cm}|}
\hline
\rowcolor{encabezado} \textbf{Verificación} & \textbf{Qué detecta} & \textbf{Resultado} \\
\hline
\archivo{tsc --noEmit} & Errores de tipos en todo el proyecto (TypeScript
estricto). & Sin errores (registro del proyecto del 3 de noviembre de 2026) \\ \hline
\archivo{eslint .} & Errores y malas prácticas, incluidas las reglas de
\emph{hooks} de React. & Sin errores (registro del 3 de noviembre de 2026) \\ \hline
\archivo{prettier --check} & Formato no uniforme del código. &
Conforme (registro del 3 de noviembre de 2026) \\
\hline
\end{tabularx}
\end{table}

\subsubsection{Pruebas unitarias automáticas}
\label{sec:estab-jest}

Como complemento de la verificación estática, se configuró Jest con el
\emph{preset} \archivo{jest-expo} y se escribieron pruebas para la lógica pura
del sistema, es decir, las funciones que no dependen de Firebase ni de módulos
nativos. Se ejecutan con \archivo{npm test}. La
Tabla~\ref{tab:suites-jest} lista los archivos de prueba de la aplicación, que
son 20 y comprenden 160 casos; en la ejecución del 3 de noviembre de 2026 todos
resultaron aprobados, igual que los 22 casos de las funciones del servidor y las
5 pruebas de JVM del módulo de ARCore (registro del proyecto en
\archivo{docs/ARCORE.md}). Su detalle se presenta en la
sección~\ref{sec:pruebas-unitarias}. \pendiente{conservar la salida de
\texttt{npm test} como evidencia para el tribunal.}

\begin{table}[H]
\centering
\caption{Suites de pruebas unitarias}
\label{tab:suites-jest}
\small
\begin{tabularx}{\textwidth}{|L{5.0cm}|Y|}
\hline
\rowcolor{encabezado} \textbf{Suite} & \textbf{Qué verifica} \\
\hline
\archivo{bmiCalculator} & Cálculo del IMC, límites de las categorías de la
OMS y mensajes. \\ \hline
\archivo{nutritionGoals} & Ecuación de Mifflin-St Jeor, factores de actividad,
ajuste por objetivo, piso calórico y coherencia de los macronutrientes. \\ \hline
\archivo{workoutStats} & Fechas, racha, semana y total de entrenamientos. \\ \hline
\archivo{workoutSession} & Creación de la sesión, series completadas y
descansos. \\ \hline
\archivo{routineGenerator} & Orden de las plantillas y esquemas (contrato con el
modelo), reglas de \archivo{decidirClases} y reparto de ejercicios. \\ \hline
\archivo{features} & Vector de entrada del modelo de rutinas y valores por
defecto. \\ \hline
\archivo{tfliteService} & Carga diferida del intérprete de TensorFlow Lite. \\ \hline
\archivo{ProgressController}, \archivo{ReportsController} & Actualización
parcial del progreso y cálculo de los reportes. \\ \hline
\archivo{MachinesController}, \archivo{machineRecognition} & Candidatas con su
caja, que no se mezclan con las de otras máquinas; preparación del objetivo de
ARCore y esquema de la respuesta del reconocimiento. \\ \hline
\archivo{workoutSessionStorage} & Persistencia de la sesión de entrenamiento
guiada. \\ \hline
\archivo{exerciseDbRepository}, \archivo{wgerImageRepository},
\archivo{localDemonstrationRepository} & Repositorios de ejercicios y
demostraciones locales por máquina. \\ \hline
\archivo{persisted} & Esquemas de los datos persistidos. \\ \hline
\archivo{environment}, \archivo{loadingProgress}, \archivo{typography} &
Variables de entorno faltantes, porcentaje de carga simulado y selección de
fuentes por peso. \\
\hline
\end{tabularx}
\fuente{Carpetas \archivo{\_\_tests\_\_} de \archivo{src/}. Las pruebas del
componente de servidor y de las reglas de seguridad se describen en la
sección~\ref{sec:pruebas-seguridad}.}
\end{table}
